% =====================================================================
% Draft section for 作業7 (paper full revision) — §5
% "Implementation Procedure" — instruction-style, to sub-subsection level.
% Compiles in BOTH arXiv (article) and TACCESS (acmart). Cites: repo,
% rp2040, arkit. Cross-refs: sec:realization, sec:clinical, sec:airback,
% tab:participants (APEE participant table in §3.3).
% Honest reproducibility note: a consolidated wiring schematic and a
% build-clean main firmware are being finalized with the device co-author;
% the repository's hardware notes track the remaining gaps.
% =====================================================================

\section{Implementation Procedure}
\label{sec:impl}

This section is written as a reproducibility guide: all software is open source
and the hardware design files (CAD, print files, firmware sketches, bill of
materials) are published in the project repository~\cite{repo}, so that a reader
can rebuild the shared iOS controller, the \emph{Pro} acoustic device (the
primary, hardware-intensive build), and the \emph{Switch} electronic device, and
re-run the APEE-style evaluation of Section~\ref{sec:clinical}. Steps reference
the realized architecture of Section~\ref{sec:realization}; paths below are
folders in the repository.

\subsection{Open-source repository and prerequisites}
\label{sec:impl-repo}

\subsubsection{Repository layout}
\label{sec:impl-layout}
The repository~\cite{repo} separates the shared controller from the two device
lines: \texttt{ios-app/} (app source plus \texttt{CODE-STRUCTURE},
\texttt{DESIGN-HIGHLIGHTS}, \texttt{BLE-CONNECTION-FLOW}, and \texttt{BUILD}
notes); \texttt{device-pro-acoustic/} (\texttt{firmware/}, \texttt{hardware/}
CAD and 3D-print files, \texttt{motor/}, \texttt{assembly/}, and a hardware
availability note); \texttt{device-switch-electronic/}; shared \texttt{docs/}
(\texttt{architecture/}, \texttt{operation/}, \texttt{toolchain/},
\texttt{user-manual/}); and \texttt{bfaaap\_patent\_info/}. Code carries no
personal signing identifiers, so a new developer supplies their own team and
bundle identifier.

\subsubsection{Toolchain prerequisites}
\label{sec:impl-toolchain}
The controller is built with Xcode and requires an Apple developer team and a
TrueDepth (Face~ID-class) iPhone or iPad, because ARKit face tracking is
device-only~\cite{arkit}. The device firmware is built from VS~Code with
PlatformIO (or the Arduino extension), the Adafruit nRF52 board support package
for the BLE board and the \texttt{arduino-pico} core for the RP2040 main
board~\cite{rp2040}. The 3D-printed parts are produced with any slicer; the
\texttt{docs/toolchain/} guide gives flashing and slicing details.

\subsection{Building the iOS controller (shared by Pro and Switch)}
\label{sec:impl-ios}

\subsubsection{Configure signing and build}
\label{sec:impl-ios-build}
Clone the repository, open \texttt{ios-app/src}, set \texttt{DEVELOPMENT\_TEAM}
and the bundle identifier to your own (placeholders are provided), keep automatic
signing, and build to a connected TrueDepth device. The \texttt{ios-app/BUILD}
note lists the exact steps; no proprietary credentials are needed.

\subsubsection{Run and verify AR$\rightarrow$BLE operation}
\label{sec:impl-ios-verify}
On first launch, grant camera permission; the play screen shows the head-tracking
state and turns red once the head passes the threshold. With a powered device
nearby, the app discovers and name-filters \texttt{bFaaaPSwitch\_n}, connects to
the selected channel, and begins streaming once the BLE write characteristic is
ready.

\subsubsection{Individualized pre-setting (offset and multiplier)}
\label{sec:impl-preset}
Before playing, the player performs a one-time \emph{individualized pre-setting}
step in the app: they choose the threshold start-point (offset) and the multiplier
to match their own comfortable range and speed of head motion, and the app stores
this preset. This step is what turns the control law of
Section~\ref{sec:control-law} into \emph{individualized control}---the same device
serves an able-footed adult, a child, and a player with restricted head motion,
each with their own offset/multiplier (e.g.\ the tracheostomy participant of
Section~\ref{sec:impact} chose a small offset and a large multiplier). Treat this
preset as a required deployment step for each new player.

\subsection{Building the Pro device (acoustic; primary)}
\label{sec:impl-pro}

\subsubsection{Mechanical assembly and 3D-printed parts}
\label{sec:impl-pro-mech}
The chassis is an aluminum-extrusion frame with 3D-printed housings (the
reference parts were printed on a 305\,mm-class printer in PLA+); the largest part
needs a bed of about $\ge$300\,mm and $\sim$200\,mm of height. The drivetrain is
motor~$\rightarrow$ timing belt $\rightarrow$ lead screw $\rightarrow$ push-rod
(Section~\ref{sec:pro-chain}). The CAD, print files, and the bill of materials
are in \texttt{device-pro-acoustic/hardware/}.

\subsubsection{Electronics, motor drive, and power}
\label{sec:impl-pro-elec}
Two boards are used: an nRF52840-class BLE board and an RP2040 main board joined
by UART. The \emph{reference} design published here uses an \emph{IQ smart-servo}
motor driven over a serial interface; because that motor is now end of life, the
\emph{next} version will use a closed-loop stepping motor (which accepts a
DRV8825-compatible STEP/DIR interface; the specific motor model is not yet
decided), but that firmware is still in development and is not the reference here
(the hardware-availability note documents the planned substitution). Pedal
reaction is sensed from \emph{motor power} rather than current, because power is
robust to supply-voltage variation (on the IQ version the power is read by command
from the motor); a commercial DC supply sized for the motor is recommended. A \emph{single-sheet KiCad
schematic} of the reference (IQ) generation --- giving the full board-to-board
wiring and the Pico pin map (BLE UART on GP0/GP1, the HX711 on GP2/GP3, the
IQ-motor serial on GP4/GP5, the air-pump MOSFET on GP12, and the slider
potentiometer on ADC0/GP26) and two power rails (+5\,V logic, +24\,V motor) ---
is included in the repository~\cite{repo}; a stepper-version schematic and a
build-clean stepper firmware are still being finalized with the device co-author.

\subsubsection{Airback (reaction anchoring) and hand controller}
\label{sec:impl-pro-air}
Wire the air pump through a power MOSFET to a controller GPIO, plumb it to an air
bag, and connect a small hand controller (an upper-limit \emph{slide
potentiometer} read on ADC0/GP26, plus a pump on/off switch). At power-up the
firmware inflates the airback until the
pressure threshold anchors the device, then self-calibrates the lower travel end
from the motor-power limit (a DIP switch sets the pressing force, $\sim$20--35\,W,
which is piano-dependent; Section~\ref{sec:airback}).

\subsubsection{Firmware flashing and device naming}
\label{sec:impl-pro-fw}
Flash the BLE board through its UF2 bootloader (Adafruit nRF52 workflow) and the
RP2040 through BOOTSEL with the \texttt{arduino-pico} core, following
\texttt{docs/toolchain/}. The advertised name (\texttt{bFaaaPSwitch\_1}..\texttt{\_4})
is stored in the BLE board's flash so several units can coexist on distinct
channels.

\subsubsection{On-piano setup and self-calibration}
\label{sec:impl-pro-setup}
Place the device on a board spanning the pedals, align the push-rod over the
sustain (damper) pedal, inflate the airback until the unit is firmly anchored,
let the actuator auto-find the lower end against the pedal, set the upper end with
the slider, and pick the on-type/off-type to match the instrument
(\texttt{docs/operation/}). The setup is non-destructive and portable, which is
what makes concert use --- not only home use --- practical.

\subsection{Building the Switch device (electronic)}
\label{sec:impl-switch}

\subsubsection{Electronics and sustain-jack interface}
\label{sec:impl-switch-elec}
The Switch needs only the nRF52840-class board plus a relay or optocoupler from a
GPIO pin to the instrument's sustain jack (a 6.3\,mm plug); there is no motor,
drivetrain, or airback. An RGB LED provides the idle/armed/engaged color
feedback.

\subsubsection{Setup}
\label{sec:impl-switch-setup}
Connect the lead to the digital piano's sustain jack, pair the channel in the
app, and set the on-type/off-type to match the jack polarity. Operation then
follows the same head-angle control law as the Pro device.

\subsection{Reproducing the APEE evaluation}
\label{sec:impl-apee}

\subsubsection{Test protocol}
\label{sec:impl-apee-protocol}
Recruit participants by class (I: able to use their own feet, as control; II:
children; III: users with leg/neck disabilities). Each participant practices the
reference score (the motif do--do--do, re--re--re, mi--mi--mi repeated four
times), selects an offset and multiplier to taste, and records three conditions:
no pedal (pattern~0), bFaaaP pedal pattern~1, and bFaaaP pedal pattern~2,
followed by a short questionnaire (the participant table is
Table~\ref{tab:participants}).

\subsubsection{Analysis and statistics}
\label{sec:impl-apee-analysis}
Import each recording into an audio analyzer (Sonic Visualiser), export the
note regions as images, and measure a tone-vibration area (TVA) with ImageJ's
particle analysis. Normalize the no-pedal condition to $1.00$ and express
patterns~1 and~2 as ratios, then compare conditions with paired $t$-tests and
compare across classes, piano experience, and instruments with one-way ANOVA.
This reproduces the result tables of Section~\ref{sec:clinical}, in which bFaaaP
significantly increased sustain ($p<0.01$) and was statistically
indistinguishable from a participant's own foot ($p>0.05$).
